\documentclass[10pt,a4paper]{amsart}
\usepackage[margin=19mm]{geometry}
\usepackage{amsmath,amssymb,mathtools}
\usepackage[scr=rsfs,cal=euler]{mathalfa}
\usepackage{xfrac}
\usepackage[unicode,hidelinks,bookmarks=false]{hyperref}
\newcommand{\ie}{i.e.\ }
\newcommand{\cf}{cf.\ }
\newcommand{\resp}{respectively\ }
\newcommand{\Subsection}{\subsection}
\providecommand{\nobreakdash}{\nobreak\hbox{-}}
\setlength{\emergencystretch}{2em}
\multlinegap0pt
\usepackage{amssymb}
\usepackage{enumerate}
\usepackage{tikz}
\newtheorem{lem}{Lemma}
\title{\LARGE B\Large\MakeLowercase{orel Exceptional Values in Several Complex Variables\newline and Their Applications to Shared Values of Shifts\newline \vspace {.25cc} and Difference Operators}}
\author{Abhijit Banerjee, Sujoy Majumder, Jhilik Banerjee}
\date{Source: arXiv:2602.14619v1 (CC BY 4.0)}
\begin{document}
\maketitle

\noindent\emph{Source and licence.} \LARGE B\Large\MakeLowercase{orel Exceptional Values in Several Complex Variables\newline and Their Applications to Shared Values of Shifts\newline \vspace {.25cc} and Difference Operators}. Version arXiv:2602.14619v1 (CC BY 4.0), distributed by arXiv under the Creative Commons Attribution 4.0 International licence, http://creativecommons.org/licenses/by/4.0/, as recorded in the arXiv licence metadata for this exact version and shown on the abstract page https://arxiv.org/abs/2602.14619v1. Full licence notice: \texttt{licenses/arxiv-2602.14619-CC-BY-4.0.txt}.\par\smallskip

\noindent\textit{Editorial scope.} Counted unit: the article's lem L2.8 (source0.tex lines 869--874). The article's complete original proof of it is delivered with it (source0.tex lines 875--892, one \texttt{proof} environment of 18 lines). No mathematics is written, repaired or re-derived: the statement and the proof are the article's own lines, in the article's own notation, and no retained line is substituted or re-flowed.  Retained uncounted as the source-local context the proof needs: lines 808--813.  Nothing was omitted from any retained span, so the package carries no omission record.  The retained lines cite 2 works, whose entries are transcribed verbatim from the article's own bibliography source in the e-print and delivered with short editorial labels recorded in the item's reference mapping.  No retained line contains a figure environment, an image-inclusion command or drawing code, so no image is delivered and the package declares no asset dependency.  Editorial formatting only, in addition: an AMS \texttt{amsart} preamble that restates the article's own declarations and macros used by the retained lines, namely the environments \texttt{lem} on separate counters, together with the standard packages those lines need.  The retained environments print in this document as Lemma 1 in order of appearance, whereas the article prints the same items with the numbering of its own full document.  The complete native source of the article as distributed in the arXiv e-print for this exact version is bundled unchanged as \texttt{sources/194730/source0.tex}.
\begin{lem}\label{L2.1}\cite[Theorem 2.4]{Cao-Xu-2020} Let $f$ be a non-constant meromorphic function on $\mathbb{C}^n$ and let $c \in \mathbb{C}^n\setminus \{0\} $. If $f$ is of finite order, then 
\begin{align*}
 m\left(r,\frac{f(z+c)}{f(z)}\right)+m\left(r, \frac{f(z)}{f(z+c)}\right)=O\left(r^{\rho(f)-1+\varepsilon}\right)
 \end{align*}
holds for any $\varepsilon>0$.
\end{lem}

\begin{lem}\label{L2.8} Let $P$ and $Q$ be two polynomials in $\mathbb{C}^n$ such that $P$ is non-constant. If $f$ is a finite order solution in $\mathbb{C}^n$ of the following equation
\begin{align*}
f(z+c)=e^{P(z)}\left(\frac{\partial f(z)}{\partial z_j}+Q(z)f(z)\right)
\end{align*}
where $c\in\mathbb{C}^n\backslash \{0\}$ and $j\in\mathbb{Z}[1,n]$, then $\rho(f)\geq \deg(P)+1$.
\end{lem}

\begin{proof} We have
\begin{align}\label{Eq3.2}
\frac{f(z+c)}{f(z)}=e^{P(z)}\left(\frac{\partial f(z)}{\partial z_j}/f(z)+Q(z)\right).
\end{align}
Now using Lemma 1.37 \cite{Hu-Li-Yang-2003}, we have
\begin{align*}
m\left(\frac{\partial f(z)}{\partial z_j}/f(z)\right)=O(\log r)
\end{align*}
and by Lemma \ref{L2.1}, we obtain
\begin{align*}
m\left(r,f(z+c)/f(z)\right)=O\left(r^{\rho(f)-1+\varepsilon}\right)
 \end{align*}
holds for any $\varepsilon>0$. Since $m(r,Q)=O(\log r)$, from (\ref{Eq3.2}), we deduce that
\begin{align*}
m(r,e^P)=O(r^{\deg(P)})\leq O\left(r^{\rho(f)-1+\varepsilon}\right)+O(\log r),
\end{align*}
which shows that $\rho(f)\geq \deg(P)+1$.
\end{proof}

\begin{thebibliography}{99}
\bibitem[Cao-Xu]{Cao-Xu-2020} {\sc T. B. Cao} and {\sc L. Xu}, Logarithmic difference lemma in several complex variables and partial difference equations, \textit{Ann. Mat. Pura Appl.}, \textbf{199} (2020), 767-794.
\bibitem[Hu-Li-]{Hu-Li-Yang-2003} {\sc P. C. Hu}, {\sc P. Li} and {\sc C. C. Yang}, Unicity of Meromorphic Mappings. Springer, New York (2003).
\end{thebibliography}
\end{document}
